% Phase 3 only: internal and entrusted manufacturing technologies.
\section{Internal and Entrusted Manufacturing Technologies}
\label{sec:p03-technologies}

This section defines the manufacturing primitives that later enter route
values. It does not construct those values, choose a route, or clear the
qualified manufacturing-service market. All technology primitives are
invariant to the MAH regime. The only direct institutional wedge remains
\(\tau_E(M)\).

\subsection{Internal manufacturing}
\label{subsec:p03-internal}

For a project with manufacturing requirement \(m>0\) and a developer with
internal manufacturing capability \(k_i>0\), let
\begin{equation}
  c_I(m,k_i)>0,
  \qquad
  c_{I,m}(m,k_i)>0,
  \qquad
  c_{I,k}(m,k_i)<0
  \quad\text{when internal production is feasible}.
  \label{eq:p03-internal-cost}
\end{equation}
The function \(c_I\) is technological marginal manufacturing cost, measured
in \(\mathsf C/\mathsf Y\). A more demanding project raises this cost, while
greater internal capability lowers it on the feasible domain.

Internal production also requires a project-level production-readiness/setup
cost \(F_I(m,k_i)\), measured in \(\mathsf C\). Let
\(\underline{k}(m)>0\) be the minimum capability needed for an internally
produced project with requirement \(m\). The feasibility convention is
\begin{equation}
  \begin{gathered}
    0\leq F_I(m,k_i)<+\infty,\quad
    F_{I,m}(m,k_i)>0,\quad F_{I,k}(m,k_i)<0
    \quad\text{if }k_i\geq\underline{k}(m),\\
    F_I(m,k_i)=+\infty
    \quad\text{if }k_i<\underline{k}(m).
  \end{gathered}
  \label{eq:p03-internal-setup}
\end{equation}
Thus low internal capability makes route \(I\) unavailable through an
infinite setup cost. No separate binary internal-feasibility primitive is
retained.

\subsection{Qualified external production}
\label{subsec:p03-external}

Under entrusted manufacturing, a qualified external producer uses the
technological marginal-cost kernel
\begin{equation}
  c_E(m)>0,
  \label{eq:p03-external-cost}
\end{equation}
also measured in \(\mathsf C/\mathsf Y\). This kernel does not depend on the
developer's internal capability \(k_i\): manufacturing is performed by the
qualified external producer.

An entrusted project requires
\begin{equation}
  b(m)>0,
  \qquad b'(m)>0,
  \label{eq:p03-capacity-requirement}
\end{equation}
units of qualified manufacturing-service capacity. The object \(b(m)\) is a
physical capacity requirement in \(\mathsf B\) per project. It is not a route
share, a supply curve, or an exogenous measure of scarcity relief.

The real technology-transfer, validation, and production-readiness cost is
\begin{equation}
  F_E(m)\geq0,
  \label{eq:p03-external-fixed-cost}
\end{equation}
measured in \(\mathsf C\) per project. In addition, the developer bears a
residual holder-side responsibility and coordination burden
\begin{equation}
  \mu_E\geq0.
  \label{eq:p03-holder-burden}
\end{equation}
The burden \(\mu_E\) is not eliminated by MAH.

The accounting roles are distinct. The technological marginal cost \(c_E(m)\)
is the input to the operating-value kernel derived in Phase~2. The later
monetary capacity payment \(p_m b(m)\), the readiness cost \(F_E(m)\), the
holder burden \(\mu_E\), and the institutional barrier \(\tau_E(M)\) are not
embedded in \(c_E(m)\). Each can therefore enter a later route value at most
once.

\subsection{Organizational and institutional distinction}
\label{subsec:p03-organization}

Routes \(I\) and \(E\) are technologically and organizationally distinct:
\begin{itemize}
  \item under \(I\), the developer remains the holder and manufactures using
  its own capability \(k_i\);
  \item under \(E\), the developer remains the authorization holder, retains
  holder responsibility, and a qualified external producer manufactures;
  \item route \(E\) is not the transfer or out-licensing route \(T\), under
  which the retained holder--producer organization is not used.
\end{itemize}

The only direct policy object is
\begin{equation}
  M\in\{0,1\},
  \qquad
  \tau_E(0)=+\infty,
  \qquad
  \tau_E(1)=\bar\tau_E<+\infty.
  \label{eq:p03-policy-invariance}
\end{equation}
The regime does not change \(c_I\), \(F_I\), \(\underline{k}\), \(c_E\),
\(b\), \(F_E\), \(\mu_E\), \(a_i\), \(k_i\), \(q\), \(m\), \(s(q)\), or
\(s_g(q)\). In particular, Phase~3 does not posit a fall in \(p_m^*\) or an
outward shift of CMO supply. Those are equilibrium objects deferred to
Phase~6.

\subsection{Local old-to-new technology map}
\label{subsec:p03-local-crosswalk}

Table~\ref{tab:p03-local-crosswalk} records only the legacy objects directly
affected by the Phase~3 definitions. The exhaustive file-, equation-,
proposition-, paragraph-, and calibration-level crosswalk remains a Phase~12
deliverable.

\begin{table}[htbp]
  \centering
  \small
  \begin{tabular}{@{}>{\raggedright\arraybackslash}p{0.23\linewidth}>{\raggedright\arraybackslash}p{0.65\linewidth}@{}}
    \hline
    Legacy object & Phase 3 treatment \\
    \hline
    \(C^I(k_i)\) &
      Replace by the distinct marginal-cost and setup-cost primitives
      \(c_I(m,k_i)\) and \(F_I(m,k_i)\). \\
    \(h_i^I\) &
      Delete as an active binary object; internal infeasibility is represented
      by \(F_I(m,k_i)=+\infty\) below \(\underline{k}(m)\). \\
    \(\bar R_i^E\) &
      Delete as a reduced-form return; Phase~2 provides \(R(q,c)\) and the
      external technology supplies \(c_E(m)\). Route-value assembly is
      deferred to Phase~4. \\
    \(q_i^E\) &
      Delete from the baseline; project requirement \(m\), qualified-external
      technology, and later capacity-market conditions handle feasibility. \\
    \(\mu_i^E\) &
      Retain the holder-side-burden concept in the common notation \(\mu_E\). \\
    \(M\) &
      Retain as the binary institutional regime. \\
    \(\eta\) &
      Remove from the baseline; it has no Phase~3 replacement. \\
    \(\zeta_i^E(M,p_m)\) &
      Remove; route \(E\) uses explicit technology, capacity requirement,
      holder burden, and the single institutional wedge. \\
    Legacy CMO scarcity shortcut &
      Do not retain as policy relief; Phase~3 defines \(b(m)\), while supply,
      demand, and \(p_m^*\) are endogenized only in Phase~6. \\
    \hline
  \end{tabular}
  \caption{Phase 3 local technology crosswalk.}
  \label{tab:p03-local-crosswalk}
\end{table}

\subsection{Boundary checks}
\label{subsec:p03-boundaries}

If \(k_i<\underline{k}(m)\), internal production is infeasible because
\(F_I=+\infty\); this is the required low-capability boundary. Increasing
\(k_i\) within the feasible domain lowers both internal marginal and setup
costs, while increasing \(m\) raises both. The entrusted technology does not
use \(k_i\), but a larger \(m\) strictly raises its qualified-capacity
requirement \(b(m)\). None of these statements is a route-choice result:
route values, outside options, and deterministic sorting are deferred to
Phase~4.
